\documentclass[12pt]{phdthesis}

\usepackage[utf8]{inputenc}
\usepackage[T1]{fontenc}
\usepackage[svgnames]{xcolor}
\usepackage{fouriernc}
\usepackage[frenchb]{babel}
% TODO \usepackage[french]{babelbib}
\usepackage{mathpartir}
\usepackage{syntax}
\usepackage{amssymb}
\usepackage{tikz}
\usepackage{amsmath}
\usepackage{float}
\usepackage{url}
\usepackage{xspace}
\usepackage{graphicx}
\usepackage[shortlabels]{enumitem}
\usepackage[backref=page]{hyperref}
\usepackage{links}
\usepackage[a4paper]{geometry}
\usepackage{amsthm}
\usepackage{stmaryrd}
\usepackage[normalem]{ulem}
\usepackage{wasysym}
\usepackage{multirow}
\usepackage{thmtools}
\usepackage{gramster}
\usepackage{switch}
\usepackage{pelagiccod}

\newcommand{\irule}[3]{ \inferrule*[right=(#1)]{#2}{#3} }

% Fix '::' -> ': :' when using babel
\newenvironment{grammarf}{\shorthandoff{:}\begin{grammar}}{\end{grammar}\shorthandon{:}}

\newcommand{\gramisa}[0]{::=~}
\newcommand{\gramor}[0]{|\hspace{3mm}}

\newcommand{\gramlr}[2]{%
  \begin{minipage}{0.2\textwidth}\raggedleft #1\end{minipage}
  \begin{minipage}{0.8\textwidth}#2\end{minipage}
}

% Types
\newcommand{\tInt}{\textsc{Int}\xspace}
\newcommand{\tFloat}{\textsc{Float}\xspace}
\newcommand{\tVoid}{\textsc{Void}\xspace}

\newcommand{\tFV}[1]{\textrm{FreeVars}(#1)}

\def\opxor{\mathbin{\char`\^}}

% Constant
\newcommand{\cNil}{\textsc{Nil}}

% Program & language names
\newcommand{\ctonewspeak}{\textsc{c2newspeak}}
\newcommand{\cpp}{\textsc{cpp}}
\newcommand{\newspeak}{\textsc{Newspeak}}

% Newspeak
\newcommand{\npkDecl}[2]{\textsc{Decl}~(#1) \{ #2 \}}
\newcommand{\npkIf}[3]{\textsc{If}(#1) \{ #2 \} \textsc{Else} \{ #3 \}}
\newcommand{\npkDoWith}[2]{\textsc{Do}\{ #1 \}\textsc{With}~#2:}
\newcommand{\npkWhile}[1]{\textsc{While}(1)\{#1\}}
\newcommand{\npkGoto}[1]{\textsc{Goto}~#1}
\newcommand{\npkReturn}[1]{\textsc{Return}~#1}

% Continuations
\def\kPass{\textsc{Kpass}}
\def\kSeq#1#2{\textsc{Kseq}(#1, #2)}
\def\kBlock#1#2{\textsc{Kblock}(#1, #2)}
\def\kGoto#1{\textsc{Kgoto}(#1)}
\def\kCall#1#2#3#4{\textsc{Kcall}(#1, #2, #3, #4)}
\def\kRet#1{\textsc{Kret}(#1)}

\newcommand{\insertcode}[1]{\input{gen/#1.pyg.tex}}

% \chapter* + add to toc
%\newcommand{\chapterstar}[1]{\chapter*{#1} \addcontentsline{toc}{chapter}{#1}}

% Register names
\def\defreg#1{\expandafter\def\csname#1\endcsname{\textsc{#1}\xspace}}

\defreg{eax}
\defreg{ebx}
\defreg{ecx}
\defreg{edx}
\defreg{esi}
\defreg{edi}
\defreg{ebp}
\defreg{esp}
\defreg{eip}
\defreg{eflags}

\def\crtrois{\textsc{CR3}}

\newcommand{\ring}[1]{\emph{ring}~#1}

\newcommand{\asminstr}[1]{\texttt{#1}}

\newcommand{\wip}[0]{%
\begin{center}
    \includegraphics[width=5cm]{gen/travaux.pdf}
\end{center}
}

\newcommand{\listcons}[0]{%
\texttt{(::)}%
}

\newcommand{\hexa}[1]{%
\texttt{0x#1}%
}

% Debug tikz bounding box
\def\drawbb{\draw[thick] (current bounding box.south west) rectangle (current bounding box.north east);}

% comes from the memoir manual {{{

\newcommand*{\addtotextheight}[1]{%
\twocolumn
\addtolength{\textheight}{#1}%
\onecolumn}

%% for main body, bottom of text at 1in, footer below
%% top of header at 1in, first text line double spaced
%% below base of header
\newlength{\linespace}
\setlength{\linespace}{\baselineskip} % current equivalent of \onelineskip
\setlength{\headheight}{\onelineskip}
\setlength{\headsep}{\linespace}
\addtolength{\headsep}{-\topskip}
\setlength{\uppermargin}{1in}
\addtolength{\uppermargin}{\headheight}
\addtolength{\uppermargin}{\headsep}

%% and for the bottom
\setlength{\lowermargin}{1in}
\setlength{\textheight}{\paperheight}

\addtolength{\textheight}{-\uppermargin}
\addtolength{\textheight}{-\lowermargin}
%% footnote settings
\setlength{\footskip}{\onelineskip}
\setlength{\footnotesep}{\onelineskip}


%% the fiddle lengths (..ta.. for title/approval page, others for prelims)
\newlength{\toptafiddle} \setlength{\toptafiddle}{2\linespace}
\newlength{\bottafiddle} \setlength{\bottafiddle}{0pt}
\newlength{\topfiddle}
\setlength{\topfiddle}{\toptafiddle}
\newlength{\botfiddle}
\setlength{\botfiddle}{\onelineskip}

%%% general macro for Abstract, etc., headings
\newcommand*{\pretoctitle}[1]{{\clearpage\centering \vspace*{-\topfiddle}#1\par}}

%%% Start the ACKNOWLEDGEMENTS
\newcommand{\asuacknowledgements}{\pretoctitle{ACKNOWLEDGEMENTS}}

%%% Start the ABSTRACT
\newcommand{\asuabstract}{ \addtotextheight{-\botfiddle} \pretoctitle{ABSTRACT}}

\addtodef{\mainmatter}{\addtotextheight{\botfiddle}{}}

%% make it easy to center any dedication
\newcommand{\asudedication}[1]{%
{\clearpage\mbox{}\vfill\centering #1 \par\vfill\clearpage}}

% }}} /memoir

\DeclareUnicodeCharacter{0393}{\Gamma}        % Γ
\DeclareUnicodeCharacter{03B5}{\epsilon}      % ε
\DeclareUnicodeCharacter{03C3}{\sigma}        % σ
\DeclareUnicodeCharacter{2190}{\leftarrow}    % ←
\DeclareUnicodeCharacter{2191}{\uparrow}      % ↑
\DeclareUnicodeCharacter{2192}{\rightarrow}   % →
\DeclareUnicodeCharacter{21D2}{\Rightarrow}   % ⇒
\DeclareUnicodeCharacter{2205}{\emptyset}     % ∅
\DeclareUnicodeCharacter{2208}{\in}           % ∈
\DeclareUnicodeCharacter{2260}{\neq}          % ≠
\DeclareUnicodeCharacter{2295}{\oplus}        % ⊕
\DeclareUnicodeCharacter{22A2}{\vdash}        % ⊢


\usetikzlibrary{calc}
\usetikzlibrary{shapes}
\usetikzlibrary{arrows}

\usepackage{fancyvrb}
\input{gen/pygments-style.tex}

\floatstyle{boxed}
\restylefloat{figure}

\addtolength{\parskip}{\baselineskip}

\title{Titre}
\author{Étienne Millon}

\begin{document}

%\maketitle

\tableofcontents

\chapterstar{Introduction}

\part{Problématique}

\chapter{Introduction}

\chapter{Rôle d'un système d'exploitation}

\section{Rôle d'un système d'exploitation}

% Fig Archi {{{
\begin{figure}
\centering
\begin{tikzpicture}
[ every node/.style={}
, node distance=2cm
]

\node (cpu) {CPU};

\node[below of=cpu] (northb) {Northbridge};

\node[right of=northb] (mem) {Mémoire};
\node[left of=northb]  (agp) {AGP};
\node[below of=northb] (southb) {Southbridge};

\node[below left of=southb, xshift=-1cm] (sata) {SATA};
\node[below right of=southb, xshift=1cm] (usb)  {USB};

\node[below left of=sata]  (hdd)     {Disque dur};
\node[below right of=sata] (cdrom)   {CD-ROM};

\node[below left of=usb]  (memo)    {Mémorette};
\node[below right of=usb] (clavier) {Clavier};

\draw (cpu) -- (northb);
\draw (northb) -- (agp);
\draw (northb) -- (mem);
\draw (southb) -- (northb);
\draw (southb) -- (sata);
\draw (southb) -- (usb);
\draw (usb) -- (memo);
\draw (usb) -- (clavier);
\draw (sata) -- (hdd);
\draw (sata) -- (cdrom);

\end{tikzpicture}

\caption{Architecture simplifiée d'un ordinateur}
\label{fig:archi-simplifiee}
\end{figure}
% }}}

Au plus bas, un ordinateur est constitué de nombreux composants matériels :
micro-processeur, mémoire, et divers périphériques
(figure~\ref{fig:archi-simplifiee}). Et au niveau de l'utilisateur, de dizaines
de logiciels permettant d'effectuer toutes sortes de calculs et de
communication. Le système d'exploitation permet de faire l'interface entre ces
niveaux d'abstraction. Dans ce chapitre nous allons voir son rôle et une
implantation possible. Pour une description plus détaillée, on pourra se référer
à l'ouvrage \cite{tanenbaum}.

Au cours de l'histoire des systèmes informatiques, la manière de les programmer
a beaucoup évolué. Au départ, les programmeurs avaient accès au matériel dans
son intégralité : toute la mémoire pouvait être accédé, toutes les instructions
pouvaient être utilisées.

Néanmoins c'est un peu restrictif : si on est seul à utiliser un système, on est
par définition... seul à pouvoir l'utiliser. Dans la seconde moitié des années
60, sont apparus les premiers systèmes "à temps partagé", permettant à plusieurs
utilisateurs de travailler en même temps.

Permettre l'exécution de plusieurs programmes en même temps est une idée
révolutionnaire, mais elle n'est pas sans difficultés techniques : en effet les
ressources de la machine doivent être aussi partagées entre les utilisateurs et
les programmes. Par exemple, plusieurs programmes vont par exemple utiliser le
CPU les uns à la suite des autres (partage \emph{temporel}) ; et chaque
programme aura à sa disposition une partie de la mémoire principale, ou du
disque dur (partage \emph{spatial}).

Si deux programmes (ou plus) s'exécutent de manière concurrente sur le même
matériel, il faut s'assurer par exemple que les deux s'exécutent à peu près
aussi souvent, ou que l'un ne puisse pas écrire dans la mémoire de l'autre. Ce
sont des rôles du système d'exploitation.

Cela passe donc par un certain bridage des possibilités du programme : plutôt
que de le faire exécuter n'importe quel type d'instruction, il communique avec
le système d'exploitation. Bien que cela ait l'air d'une limitation, c'est aussi
bénéfique pour le programmeur puisque cela permet de définir des abstractions au
niveau du noyau.

Par exemple, si un programmeur veut copier des données depuis un CD-ROM vers la
mémoire principale, il devra interroger le bus SATA, interroger le lecteur sur
la présence d'un disque dans le lecteur, activer le moteur, calculer le numéro
de \emph{frame} des données sur le disque, demander la lecture, puis déclencher
une copie de la mémoire.

Si dans un autre cas il voulait récupérer des données depuis une mémorette USB,
il devrait interroger le bus USB, rechercher le bon numéro de périphérique, le
bon \emph{endpoint} dans celui-ci, lui appliquer une commande de lecture au bon
numéro de bloc, puis copier la mémoire.

Ces deux opérations, bien qu'elles aient le même but (copier de la mémoire
depuis un périphérique amovible), ne sont pas effectuées en pratique de la même
manière. C'est pourquoi le système d'exploitation fournit les notions de
fichier, lecteur, etc : le programmeur n'a plus qu'à utiliser des commandes de
haut niveau (``monter un lecteur'', ``ouvrir un fichier'', ``lire dans un
fichier'') et selon le type de lecteur, le système d'exploitation effectuera les
actions appropriées.

En résumé, un système d'exploitation est l'intermédiaire entre le logiciel et
le matériel, et en particulier assure les rôles suivants :

\paragraph{Gestion des processus} : un système d'exploitation peut permettre
d'exécuter plusieurs programmes à la fois. Il faut alors orchestrer ces
différents processus et les séparer en terme de temps et de ressources
partagées.

\paragraph{Gestion de la mémoire} : chaque processus, en plus du noyau, doit
disposer d'un espace mémoire différent. C'est-à-dire qu'un processus ne doit pas
pouvoir interférer avec un autre.

\paragraph{Gestion des périphériques} : le noyau étant le seul code à s'exécuter
en mode privilégié, c'est lui qui doit communiquer avec les périphériques
matériels.

\paragraph{Abstractions} : le noyau fournit aux programmes une interface
unifiée, permettant de stocker des informations de la même manière sur un disque
dur ou une clef USB (alors que l'accès se déroulera de manière très différente
en pratique). C'est ici que la notion arbitraire de fichier sera définie, par
exemple.

\section{Architecture Intel}

Dans cette section nous décrivons le fonctionnement d'un processeur utilisant
une architecture Intel 32 bits. Les morceaux de code seront écrit en syntaxe
AT\&T, celle que comprend l'assembleur GNU.

\subsection{Assembleur}

Pour faire des calculs, le processeur est composé de registres, qui sont des
petites zones de mémoire interne, et peut accéder à la mémoire principale.

Les registres principaux sont nommés \eax, \ebx, \ecx, \edx, \esi et \edi. Ils
peuvent être utilisés pour n'importe quel type d'opération, mais certains sont
spécialisés : par exemple il est plus efficace d'utiliser \eax en accumulateur,
ou \ecx en compteur. \esp est le pointeur de pile (\emph{stack pointer}), qui a
un rôle particulier dans les instructions \asminstr{push} et \asminstr{pop}.
\ebp sert de point de repère dans les appels de fonction. \eip contient
l'adresse de l'instruction courante.

Les calculs sont décrits sous forme d'une suite d'instructions. Chaque
instruction est composée d'une mnémonique et d'une liste d'opérandes. Les
mnénoniques (\asminstr{mov}, \asminstr{call}, \asminstr{sub}, etc) définissent
un type d'opération à appliquer sur les opérandes.

Ces opérandes peuvent être de plusieurs types :

\begin{itemize}
\item
  un nombre, noté \$4
\item
  le nom d'un registre, noté \texttt{\%eax}
\item
  une opérande mémoire, c'est à dire le contenu de la mémoire à une
  adresse effective. Cette adresse effective peut être exprimée de
  plusieurs manières :
  \begin{itemize}
  \item
    directement : \texttt{addr}
  \item
    indirectement : \texttt{(\%ecx)}. L'adresse effective est le contenu
    du registre.
  \item
    ``base + déplacement'' : \texttt{4(\%ecx)}. L'adresse effective est
    le contenu du registre plus le déplacement (4 ici).
  \end{itemize}
\end{itemize}

En pratique il y a des modes d'adressage plus complexes, et toutes les
combinaisons ne sont pas possibles, mais ceux-ci suffiront à décrire les
exemples suivants :

\begin{itemize}

\item \texttt{mov src, dst} copie le contenu de \texttt{src} dans \texttt{dst}.

\item \texttt{add src, dst} calcule la somme des contenus de \texttt{src} et
  \texttt{dst} et place ce résultat dans \texttt{dst}.

\item \texttt{push src} place \texttt{src} sur la pile, c'est à dire que cette
  instruction décrémente le pointeur de pile \esp de la taille de \texttt{src},
  puis place \texttt{src} à l'adresse mémoire de la nouvelle valeur \esp.

\item \texttt{pop src} réalise l'opération inverse : elle charge le contenu de
  la mémoire à l'adresse \esp dans \texttt{src} puis incrémente \esp de la
  taille correspondante.

\item \texttt{jmp addr} saute à l'adresse \texttt{addr} : c'est l'équivalent de
  \texttt{mov addr, \%eip}.

\item \texttt{call addr} sert aux appels de fonction : cela revient à
  \texttt{push \%eip} puis \texttt{jmp addr}.

\item \texttt{rest} sert à revenir d'une fonction : c'est l'équivalent de
  \texttt{pop \%eip}.

\end{itemize}

\subsection{Appel de fonction et conventions d'appel}

\begin{figure} % {{{ fig:stackframe
\begin{tikzpicture}
[ stack/.style={draw,shape=rectangle,minimum height=1cm,minimum width=5cm}
, bigbrace/.style={decorate,decoration={brace,amplitude=3mm}}
, >=latex
]

  \path (0,0) node [stack] (gloc) {Locales de g}
  -- ++(0,-1) node [stack] (gebp) {}
  -- ++(0,-1) node [stack] (gret) {}
  -- ++(0,-1) node [stack] (gpar) {Arguments de g}
  -- ++(0,-1) node [stack] (floc) {Locales de f}
  -- ++(0,-1) node [stack] (febp) {}
  -- ++(0,-1) node [stack] (fret) {}
  -- ++(0,-1) node [stack] (fpar) {Arguments de f};

  \draw [->] (gebp) -- ++(4, 0) |- (febp.5);
  \draw [->] (febp.355) -- ++(1.5, 0) -- ++(0,-4);

  \draw [dashed] (fpar.south west) -- ++ (0,-2);
  \draw [dashed] (fpar.south east) -- ++ (0,-2);
  \draw [dashed] (gloc.north east) -- ++ (0,2);
  \draw [dashed] (gloc.north west) -- ++ (0,2);

  \node at (5, 1) (ox) {};

  \node [text width=3cm] at (-0.8, 2)    {Adresses hautes (sens de \asminstr{push})};
  \node [text width=3cm] at (-0.8, -8.5) {Adresses basses (sens de \asminstr{pop})};

  \node [text width=3cm] at (-6, -3) (asmtext) {\tt
    g:  ... \\
    leave \\
    ret \\
    f:  ... \\
    pushl \$2 \\
    pushl \$4 \\
    call g \\
    addl \$8, \%esp \\
    leave \\
    ret \\
    ... \\
    call f\\
    ...
    };

  \draw [-*] (gret) -- ++ (-3, 0) |- ($ (asmtext.north west) + (2, -3.5) $);
  \draw [-*] (fret) -- ++ (-3, 0) |- ($ (asmtext.north west) + (2, -6) $);

  \draw [<-] (gebp.west) -- ++ (-0.5,0) node [xshift=-4mm] {\ebp};
  \draw [<-] (gloc.west) -- ++ (-0.5,0) node [xshift=-4mm] {\esp};

  \draw [<-] ($ (asmtext.north west) + (2, 0) $) -- ++ (0.5,0) node [xshift=3mm] {\eip};

  \draw [bigbrace]
        ($ (gloc.north east) + (3, 0) $)
     -- node [auto,midway,xshift=3mm] {Cadre de g}
        ($ (gpar.south east) + (3, 0.1) $);

  \draw [bigbrace]
        ($ (floc.north east) + (3, -0.1) $)
     -- node [auto,midway,xshift=3mm] {Cadre de f}
        ($ (fpar.south east) + (3, 0) $);

\end{tikzpicture}

\caption[Cadres de pile]{ Cadres de pile. }
\label{fig:stackframe}
\end{figure} % }}}

Dans le langage d'assemblage, il n'y a pas de notion de fonction ; mais
\asminstr{call} et \asminstr{ret} permettent de sauvegarder et de restaurer une
adresse de retour, ce qui permet de faire un saut et revenir à l'adresse initiale.
Ce système permet déjà de créer des procédures, c'est-à-dire des fonctions sans
arguments ni valeur de retour.

Pour gérer ceux-ci, il faut que les deux morceaux (appelant et appelé) se
mettent d'accord sur une convention d'appel commune. La convention utilisée sous
GNU/Linux est appelée \emph{cdecl} et possède les caractéristiques suivantes :

\begin{itemize}
\item la valeur de retour d'une fonction est stockée dans \eax
\item \eax, \ecx et \edx peuvent être écrasés sans avoir à les sauvegarder
\item les arguments sont placés sur la pile (et enlevés) par l'appelant. Les
  paramètres sont empilés de droite à gauche.
\end{itemize}

Pour accéder à ses paramètres, une fonction peut donc utiliser un adressage
relatif à \esp. Cela peut fonctionner, mais cela complique les choses si elle
contient aussi des variables locales. En effet, les variables locales sont
placées sur la pile, au dessus des (c'est à dire, empilées après) paramètres,
augmentant le décalage.

Pour simplifier, la pile est organisée en cadres logiques : chaque cadre
correspond à un niveau dans la pile d'appels de fonctions. Si f appelle g, qui
appelle h, il y aura dans l'ordre sur la pile le cadre de f, celui de g puis
celui de h.

Ces cadres sont chainés à l'aide du registre \ebp : à tout moment, \ebp contient
l'adresse du cadre de l'appelant.

Prenons exemple sur la figure~\ref{fig:stackframe} : pour appeler
\texttt{g(4,2)}, \texttt{f} empile les arguments de droite à gauche.
L'instruction \asminstr{call g} empile l'adresse de l'instruction suivante sur
la pile puis saute au début de \texttt{g}.

Au début de la fonction, les trois instructions permettent de sauvegarder
l'ancienne valeur de \ebp, faire pointer \ebp à une endroit fixe dans le cadre
de pile, puis allouer 8 octets de mémoire pour les variables locales.

Dans le corps de la fonction \texttt{g}, on peut donc se référer aux variables
locales par \texttt{-4(\%ebp)}, \texttt{-8(\%ebp)}, etc, et aux arguments par
  \texttt{8(\%ebp)}, \texttt{12(\%ebp)}, etc.

À la fin de la fonction, l'instruction \asminstr{leave} est équivalente à
\texttt{mov \%ebp, \%esp} puis \texttt{pop \%ebp} et permet de défaire le cadre
de pile, laissant l'adresse de retour en haut de pile. Le \asminstr{ret} final
la dépile et y saute.

\subsection{Tâches, niveaux de privilèges}

Sans mécanisme particulier, le processeur exécuterait uniquement une suite
d'instruction à la fois. Pour lui permettre d'exécuter plusieurs tâches, un
système de partage du temps existe.

À des intervalles de temps réguliers, le système est programmé pour recevoir une
interruption. C'est une condition exceptionnelle (au même titre qu'une division
par zéro) qui fait sauter automatiquement le processeur dans une routine de
traitement d'interruption. Si à cet endroit le code sauvegarde les registres et
restaure un autre ensemble de registres, on peut exécuter plusieurs tâches de
manière entrelacée. Si l'alternance est assez rapide, cela peut donner
l'illusion que les programmes s'exécutent en parallèle. Comme l'interruption
peut survenir à tout moment, on parle de multitâche préemptif.

En plus de cet ordonnancement de processus, l'architecture Intel permet
d'affecter des niveaux de privilège à ces tâches, en restreignant le type
d'instructions exécutables, ou en donnant un accès limité à la mémoire aux
tâches de niveaux moins elevés.

Il y a 4 niveaux de privilèges (nommés aussi \emph{rings}) : le \ring{0} est le
plus privilégié, le \ring{3} le moins privilégié. Dans l'exemple précédent, on
pourrait isoler l'ordonnanceur de processus en le faisant s'exécuter en \ring{0}
alors que les autres tâches seraient en \ring{3}.

\begin{center}\rule{3in}{0.4pt}\end{center}

\begin{figure}
\begin{tikzpicture}

\begin{scope}

  \path[clip] (-1, 0.8) rectangle (5, -0.8);

  \foreach \x/\c in { 4/20, 3/30, 2/40, 1/50, 0/60 }
  \draw[fill=black!\c] ($ (-2,0) + (\x,0) $) circle (2 cm);

\end{scope}

%\draw [<-] (3.5, 0) -- ++(0,-2) node {Applications};
%\draw [<-] (0, 0) -- ++(0,-2) node {Noyau};

\end{tikzpicture}

\caption[Les différents \emph{rings}]{
  Les différents \emph{rings}. Seul le \ring{0} a accès au hardware
  via des instructions privilégiées. Pour accéder aux fonctionnalités du noyau,
  les programmes utilisateur doivent passer par des appels système.}
\label{fig:rings}
\end{figure}

La première protection nécessaire est d'isoler le système d'exploitation
lui-même des programmes qu'il va permettre d'exécuter. Pour ce faire, Le
processeur permet d'exécuter des tâches selon plusieurs niveaux de privilège,
aussi appelés \emph{rings} (figure~\ref{fig:rings}) : du \ring{3}, le
moins privilégié, jusqu'au \ring{0}, le plus privilégié. On peut configurer
le processeur de manière à ce que les instructions privilégiées (accès aux ports
d'entrée/sortie...) ne soient possibles qu'en \ring{0}. Bien que 4 niveaux
soient disponibles, Linux n'utilise que les \emph{rings} 0 et 3 : le noyau
lui-même en \ring{0} et les processus utilisateur en \ring{3}.

\begin{figure} % fig:memoire-virtuelle {{{
\centering
\begin{tikzpicture}
  [every node/.style={node distance=2cm}]

  \node (u1) {};

  \node[right of=u1] (phy) {};

  \node[right of=phy] (u2) {};

  \draw (u1) rectangle ++(1,-6);
  \draw (phy) rectangle ++(1,-6);
  \draw (u2) rectangle ++(1,-6);
\end{tikzpicture}
\caption{Mécanisme de mémoire virtuelle.}
\label{fig:memoire-virtuelle}
\end{figure} % }}}

En ce qui concerne la mémoire, les différentes tâches ont une vision différente
de la mémoire physique : c'est-à-dire que deux tâches lisant à une même adresse
peuvent avoir un résultat différent. C'est le concept de mémoire virtuelle
(fig~\ref{fig:memoire-virtuelle}).

\begin{figure} % fig:memmap {{{
\centering
\fbox{
  \begin{tikzpicture}
  [scale=0.7
  ,user/.style={fill=black!20}
  ,kernel/.style={fill=black!70}
  ]

  % Memory zone
  %
  % #1 - start
  % #2 - end
  % #3 - color
  \newcommand{\mzone}[3]{
    \path[#3] (#1,0) rectangle (#2,1);
  }

  % Address label
  %
  % #1 - x position
  % #2 - text
  \newcommand{\alabel}[2]{
    \path (#1,1) -- ++(0,0.3) node [pos=1] {\small \ttfamily #2};

  }

  % exec
  \mzone{0.5}{1}{user}

  % lib
  \mzone{2.5}{3.2}{user}

  % stack
  \mzone{3.7}{4}{user}

  % stack
  \mzone{5}{5.5}{user}

  % kernel
  \mzone{6}{8}{kernel}

  % contour
  \draw (0,0) rectangle (8,1);

  \alabel{0}{0}
  \alabel{6}{3 Go}
  \alabel{8}{4 Go}

\end{tikzpicture}

}

\caption[Espace d'adressage d'un processus]{L'espace d'adressage d'un processus.
En gris clair, les zones accessibles à tous les niveaux de privilèges : code du
programme, bibliothèques, tas, pile. En gris foncé, la mémoire du noyau,
réservée au mode privilégié.}

\label{fig:memmap}
\end{figure}
% }}}

\section{Cas de Linux}

Dans cette section, nous allons voir comment ses mécanismes sont implantés dans
le noyau Linux. Une description plus détaillée pourra être trouvée dans
\cite{UnderstandingTheLinuxKernel} par exemple.


\subsection{Appels système}

Les programmes utilisateur s'exécutant en \ring{3}, ils ne peuvent pas
contenir d'instructions privilégiées, et donc ne peuvent pas accéder directement
au matériel (c'était le but !). Pour qu'ils puissent interagir avec le système
(afficher une sortie, écrire sur le disque...), le mécanisme des appels système
est nécessaire. Il s'agit d'une interface de haut niveau entre les \emph{rings}
3 et 0. Du point de vue du programmeur, il s'agit d'un ensemble de fonctions C
``magiques'' qui font appel au système d'exploitation pour effectuer des
opérations.

Prenons le cas de l'appel \texttt{getpid}, qui retourne le numéro de processus
courant. La bibliothèque C fournit une fonction du même nom :

\begin{Verbatim}
pid_t getpid(void);
\end{Verbatim}

A priori, rien de différent d'une fonction implantée directement en C. Par un
processus détaillé ci-après, cette fonction va invoquer la fonction, suivante,
définie dans le noyau (\texttt{kernel/timer.c}) :

\begin{Verbatim}
SYSCALL_DEFINE0(getpid)
{
        return task_tgid_vnr(current);
}
\end{Verbatim}



Le mécanisme de couplage entre ces deux fonctions est le suivant. Une
explication plus détaillé est disponible dans la documentation d'Intel
\cite{intelsys}.

Il y a bien une fonction \texttt{getpid} présente dans la bibliothèque C du
système. C'est la fonction qui est directement appelée par le programme. Cette
fonction commence par placer le numéro de l'appel système (noté
\texttt{\_\_NR\_getpid}, valant 20 ici) dans \eax, puis les arguments éventuels
dans les registres (\ebx, \ecx, \edx, \esi puis \edi). Une interruption
logicielle est ensuite déclenchée (\texttt{int 0x80}) : étant donné la
configuration du processeur, elle sera traitée en \ring{0}, à un point
d'entrée prédéfini (\texttt{arch/x86/kernel/entry\_32.S},
\texttt{ENTRY(system\_call)}).

L'exécution reprend donc en \ring{0}, avec une pile propre au processus.
Les valeurs des registres ont été préservées : elles sont mises sur la pile du
noyau. Le noyau examine le numéro d'appel système (toujours dans \eax) et
appelle la fonction correspondante dans le tableau \texttt{sys\_call\_table} :
les arguments sont en place sur la pile, comme dans le cas d'un appel de
fonction ``classique''. La fonction, nommée \texttt{sys\_getpid}, peut donc être
écrite en C.

Au retour de la fonction, la valeur de retour est placée à la place de \eax là
où les registres ont été sauvegardés sur la pile noyau. L'instruction
\texttt{iret} permet de restaurer les registres et de repasser en mode
utilisateur, juste après l'interruption. La fonction de la bibliothèque C peut
alors retourner au programme appelant.

% vim: spelllang=fr


\chapter{Cas d'étude}

On a vu que les appels systèmes permettent aux programmes utilisateur d'accéder
au services du noyau. Ils forment donc une interface particulièrement sensible aux
problèmes de sécurité.

Comme pour toutes les interfaces, on peut être plus ou moins fin. D'un
côté, une interface pas assez fine serait trop restrictive et ne permettrait pas
d'implémenter tout type de logiciel. De l'autre, une interface trop laxiste
("écrire dans tel registre matériel") empêche toute isolation. Il faut donc
trouver la bonne granularité.

Nous allons présenter ici une difficulté liée à la manipulation de mémoire au
sein de certains types d'appels système.

Il y a deux grands types d'appels systèmes : d'une part, ceux qui renvoient un
simple nombre, comme \texttt{getpid} qui renvoie le numéro du processus
appelant.

\begin{Verbatim}
int main(void)
{
    pid_t pid = getpid();
    printf("%d\n", pid);
}
\end{Verbatim}

Ici, pas de difficulté particulière : la communication entre le \emph{ring} 0 et le
\emph{ring} 3 est faite uniquement à travers les registres, comme décrit dans le
chapitre précédent.

\begin{center}\rule{3in}{0.4pt}\end{center}





Dans le cas d'un appel système qui ne fait que renvoyer un nombre, il n'y a pas
de difficulté. En revanche, certains appels système (la majorité, en fait),
remplissent une structure avec leurs résultats :

\begin{Verbatim}
int gettimeofday(struct timeval *tv, struct timezone *tz);
\end{Verbatim}

Pour utiliser une telle fonction, l'utilisateur doit allouer lui même la
mémoire, sur la pile par exemple :

\begin{Verbatim}
struct timeval tv;
struct timezone tz;
int z = gettimeofday(&tv, &tz);
if (z == 0) {
    printf( "tv.tv_sec = %ld\ntv.tv_usec = %ld\n"
            "tz.tz_minuteswest = %d\ntz.tz_dsttime = %d\n",
             tv.tv_sec, tv.tv_usec,
             tz.tz_minuteswest, tz.tz_dsttime
          );
}
\end{Verbatim}

Notons que dans ce cas, c'est le noyau qui remplit la structure : le
déréférencement se fait en \emph{ring} 0.

\begin{center}\rule{3in}{0.4pt}\end{center}



\section{Description du problème}

% But
% What you're trying to tell the audience that they don't already know (e.g. Your story.)

Un système d'exploitation moderne comme GNU/Linux est séparé en deux niveaux de
privilèges : le noyau, qui gère directement le matériel, et les applications de
l'utilisateur, qui communiquent avec le noyau par l'interface restreinte des
\emph{appels système}.

Pour assurer l'isolation, ces deux parties n'ont pas accès aux mêmes zones
mémoire (cf. figure~\ref{fig:memmap}).

Si le code utilisateur tente d'accéder à la mémoire du noyau, une erreur sera
déclenchée. En revanche, si cette écriture est faite au sein de l'implantation
d'un appel système, il n'y aura pas d'erreur puisque le noyau a accès à toute la
mémoire : l'isolation aura donc été brisée.

Pour celui qui implante un appel système, il faut donc empêcher qu'un pointeur
passé en paramètre référence le noyau. Autrement dit, il est indispensable de
vérifier dynamiquement que la zone dans laquelle pointe le paramètre est
accessible par l'appelant\cite{hardy88confused}.

Si au contraire un tel pointeur est déréférencé sans vérification (avec
\texttt{*} ou une fonction comme \texttt{memcpy}), le code s'exécutera
correctement mais en rendant le système vulnérable, comme le montre la
figure~\ref{fig:radeon-bug}.

% Figure bug {{{
\begin{figure}
  \insertcode{radeon-buggy.c}

  \caption{Bug freedesktop.org \#29340. Le paramètre \texttt{data} provient de
    l'espace utilisateur via un appel système. Un appelant malveillant peut se
    servir de cette fonction pour lire la mémoire du noyau à travers le message
    d'erreur.}

  \label{fig:radeon-bug}
\end{figure}
% }}}

Pour éviter cela, le noyau fournit un ensemble de fonctions qui permettent de
vérifier dynamiquement la valeur d'un pointeur avant de le déréférencer. Par
exemple, dans la figure précédente, la ligne 8 aurait dû être remplacée par :

\begin{Verbatim}
copy_from_user(&value, value_ptr, sizeof(value));
\end{Verbatim}

L'analyse présentée ici permet de vérifier automatiquement et statiquement que
les pointeurs qui proviennent de l'espace utilisateur ne sont déréférencés qu'à
travers une de ces fonctions sûres.

\section{Principes de l'analyse}

% Why the audience should believe that the results you've got aren't made up or
% flawed
Le problème est modélisé de la façon suivante : on associe à chaque variable
\texttt{x} un type de données \texttt{t}, ce que l'on note \texttt{x:t}. En
plus des types présents dans le langage C, on ajoute une distinction
supplémentaire pour les pointeurs. D'une part, les pointeurs ``noyau'' (de type
\texttt{t~*}) sont créés en prenant l'adresse d'un objet présent dans le code
source. D'autre part, les pointeurs ``utilisateurs'' (leur type est noté
\texttt{t user*}) proviennent des interfaces avec l'espace utilisateur.

Il est sûr de déréférencer un pointeur noyau, mais pas un pointeur
utilisateur. L'opérateur \texttt{*} prend donc un \texttt{t *} en entrée
et produit un \texttt{t}.

Pour faire la vérification de type sur le code du programme, on a besoin de
quelques règles. Tout d'abord, les types suivent le flot de données.
C'est-à-dire que si on trouve dans le code \texttt{a = b}, \texttt{a} et
\texttt{b} doivent avoir un type compatible. Ensuite, le qualificateur
\texttt{user} est récursif : si on a un pointeur utilisateur sur une structure,
tous les champs pointeurs de la structure sont également utilisateur. Enfin, le
déréférencement s'applique aux pointeurs noyau seulement : si le code contient
l'expression \texttt{*x}, alors il existe un type \texttt{t} tel que
\texttt{x:t*} et \texttt{*x:t}.

Appliquons ces règles à l'exemple de la figure \ref{fig:radeon-bug} : on suppose
que l'interface avec l'espace utilisateur a été correctement annotée. Cela
permet de déduire que \texttt{data:void user*}. En appliquant la première règle
à la ligne 6, on en déduit que \texttt{info:struct drm\_radeon\_info user*}
(comme en C, on peut toujours convertir de et vers un pointeur sur
\texttt{void}).

Pour déduire le type de \texttt{value\_ptr} dans la ligne 7, c'est la
deuxième règle qu'il faut appliquer : le champ \texttt{value} de
la structure est de type \texttt{uint32\_t~*} mais on y accède à travers
un pointeur utilisateur, donc \texttt{value\_ptr:uint32\_t user*}.

À la ligne 8, on peut appliquer la troisième règle : à cause du déréférencement,
on en déduit que \texttt{value\_ptr:t *}, ce qui est une contradiction puisque
d'après les lignes précédentes, \texttt{value\_ptr:uint32\_t user*}.

Si la ligne 3 était remplacée par l'appel à \texttt{copy\_from\_user}, il n'y
aurait pas d'erreur de typage car cette fonction peut accepter les arguments
\texttt{(uint32\_t~*, uint32\_t user*, size\_t)}.

\section{Implantation}

Une implantation est en cours. Le code source est d'abord prétraité par
\texttt{gcc -E} puis converti en Newspeak~\cite{newspeak}, un langage destiné à
l'analyse statique. Ce traducteur peut prendre en entrée tout le langage C, y
compris de nombreuses extensions GNU utilisées dans le noyau. En particulier,
l'exemple de la figure~\ref{fig:radeon-bug} peut être analysé.

À partir de cette représentation du programme et d'un ensemble d'annotations
globales, on propage les types dans les sous-expressions jusqu'aux feuilles.

Si aucune contradiction n'est trouvée, c'est que le code respecte la propriété
d'isolation. Sinon, cela peut signifier que le code n'est pas correct, ou bien
que le système de types n'est pas assez expressif pour le code en question.

Le prototype, disponible sur \url{http://penjili.org}, fera l'objet d'une
démonstration.

\section{Conclusion}

% Recap of your story and its implications

Nous avons montré que le problème de la manipulation de pointeurs non sûrs peut
être traité avec une technique de typage. Elle est proche des analyses menées
dans CQual~\cite{pldi99} ou Sparse~\cite{TorvaldsSparse}.

% Limitations
% Why someone might doubt your story and what you've done to get rid of as much
% doubt as possible.

Plusieurs limitations sont inhérentes à cette approche : notamment, la présence
d'unions ou de \emph{casts} entre entiers et pointeurs fait échouer l'analyse.

% Extensions

Le principe de cette technique (associer des types aux valeurs puis restreindre
les opérations sur certains types) peut être repris. Par exemple, si on définit
un type ``numéro de bloc'' comme étant un nouvel alias de \texttt{int}, on peut
considérer que multiplier deux telles valeurs est une erreur.


\chapter{État de l'art}

\begin{itemize}
\item
  Analyse légère :
\item
  Sparse \cite{TorvaldsSparse}
\item
  Analyse de valeurs :
\item
  IA :
  \begin{itemize}
  \item
    Cousot 77 \cite{Cousot77}
  \item
    Cousot 92 \cite{Cousot92-1}
  \end{itemize}
\item
  Typage fort
\item
  \cite{TAPL}
\item
  Qualifiers in Haskell :
  \begin{itemize}
  \item
    \cite{lightweight-static-capabilities}
  \item
    \cite{LZ06a}
  \end{itemize}
\item
  Qualificateurs de type :
\item
  CQual \cite{pldi99,usenix01,pldi02,cquk-usenix04,toplas-quals}
\end{itemize}


\part{Analyse statique par typage}

\chapter{Typage}

\section{Présentation et but}

Au plus bas niveau d'abstraction, un ordinateur ne manipule que des nombres
entiers \footnote{en allant plus loin on pourrait dire qu'il ne manipule que des
suites de bits} : en langage machine, il n'y a pas de distinction entre une
adresse et un nombre.

Pourtant il est clair que certaines opérations n'ont pas de sens : par exemple,
ajouter deux adresses, ou déréférencer le résultat d'une division sont des
comportements qu'on voudrait pouvoir empêcher.

En un mot, le but du typage est de classifier les objets et de restreindre les
opérations possibles selon la classe d'un objet : "ne pas ajouter des pommes et
des oranges".

Le modèle qui permet cette classification est appelé \emph{système de types} et
est en général constitué d'un ensemble de \emph{règles de typage}, comme "un
entier plus un entier égale un entier".

\section{Taxonomie}

La définition d'un langage de programmation introduit la plupart du temps celle
d'un système de types. Il y a donc de nombreux systèmes de types différents,
dont nous pouvons donner une classification sommaire.

\subsection{Dynamique, statique, mixte}

Il y a deux grandes familles de systèmes de types, selon quand se fait la
vérification de types. On peut en effet l'effectuer au moment de l'exécution, ou
au contraire prévenir les erreurs à l'exécution en la faisant au moment de la
compilation (ou avant l'interprétation).

\subsubsection{Typage dynamique}

La première est le typage \emph{dynamique}. Pour différencier les différents
types de données, on ajoute une étiquette à chaque valeur. Dans tout le
programme, on ne manipulera que des valeurs étiquettées :

\begin{itemize}
\item
  si on veut réaliser l'opération
  $(0x00000001, Int) + (0x0000f000, Int)$, on vérifie tout d'abord qu'on
  peut réaliser l'opération $+$ entre deux $Int$. Ensuite on réalise
  l'opération elle même, qu'on étiquette avec le type du résultat :
  $(0x0000f001, Int)$
\item
  si au contraire on tente d'ajouter deux adresses
  $(0x2e8d5a90, Addr) + (0x76a5e0ec, Addr)$, la vérification échoue et
  l'opération s'arrête avec une erreur.
\end{itemize}

Il existe plusieurs techniques pour signaler les erreurs de typage dynamiques :
arrêter l'exécution, lever une exception, convertir une opérande, utiliser une
valeur d'erreur, etc.

\subsubsection{Typage statique}

La seconde technique est le typage \emph{statique} : plutôt que de vérifier les
types sur les données, on les vérifie à l'arrêt, sur les expressions. Cela
implique par exemple que chaque variable doit contenir des valeurs d'un même
type tout au long de sa vie.

\todo{...}

\subsubsection{Typage hybride ("stanamique")}

Il est aussi possible de mélanger ces deux approches :

\begin{itemize}
\item
  à la compilation, essayer d'obtenir les types les plus précis
  possibles
\item
  pour les cas restants, insérer un test dynamique
\end{itemize}

\subsection{Fort, faible, sound}

Un autre critère intéressant est la "force" du typage : est-ce qu'il peut
induire des conversions de type implicites ?

Prenons l'exemple de l'addition entre un entier $2$ et un flottant $3.14$. On
peut le voir de plusieurs manières différentes :

\begin{itemize}
\item
  un entier est un nombre décimal comme les autres : $2 = 2.00$. Donc
  $2 + 3.14 = 2.00 + 3.14 = 5.14$.
\item
  $2$ est un entier, $3.14$ un flottant : ces deux types sont différents et
  l'opération est impossible.
\end{itemize}

\subsection{Polymorphisme}

Dans le cas du typage statique, restreindre une opération à un seul type de
données peut être assez restrictif.

Par exemple, quel doit être le type d'une fonction qui trie un tableau en place
?

\begin{figure}

\begin{tikzpicture}
  [ ptype/.style={draw=black,shape=rectangle}
  , child/.style={ptype, xshift=2cm}
  , gen0/.style={       node distance=2cm}
  , gen1/.style={child, node distance=1cm}
  , gen2/.style={child, node distance=5mm}
  ]

  \node[ptype]                     (mono) {Mono};
  \node[ptype,below of=mono, gen0] (poly) {Poly};
  \node[above right of=poly, gen1] (univ)  {Universel};
  \node[above right of=univ, gen2] (param) {Paramétrique};
  \node[below right of=univ, gen2] (inclu) {Inclusion};
  \node[below right of=poly, gen1] (adho)  {Ad-hoc};
  \node[above right of=adho, gen2] (over)  {Surcharge};
  \node[below right of=adho, gen2] (coer)  {Coercition};

  \draw (poly) -- (univ);
  \draw (poly) -- (adho);

  \draw (univ) -- (param);
  \draw (univ) -- (inclu);

  \draw (adho) -- (over);
  \draw (adho) -- (coer);

      %node {Mono}
      %node {Poly}
          %node {Universal}
              %node {Param}
              %node {Inclusion}
          %node {Ad-Hoc}
              %node {Surcharge}
              %node {Coercition}

\end{tikzpicture}

\caption{Les différents types de polymorphisme.}
\label{fig:types-de-polymorphisme}
\end{figure}

\subsubsection{Monomorphisme}

Une première solution peut être de forcer des types concrets, c'est à dire
qu'une même fonction ne pourra s'appliquer qu'à un seul type de données.

\subsubsection{Polymorphisme paramétrique}

\subsubsection{Polymorphisme par inclusion}

\subsubsection{Polymorphisme par surcharge}

\subsubsection{Polymorphisme par coercition}

\subsection{Expressivité, garanties, types dépendants}

\section{Exemples}

\subsection{Faible dynamique : Perl}
\subsection{Faible statique : C}
\subsection{Fort dynamique : Python}
\subsection{Fort statique : OCaml}
\subsection{Fort statique à effets typés : Haskell}
\subsection{Theorem prover : Coq}


\chapter{Un premier système de types}

\input{chap/typbase.tex}

\chapter{Système avec le qualificateur "user"}

\input{chap/qualuser.tex}

\chapter{Système avec le qualificateur "sz"}

\chapter{Implantation}

\chapterstar{Conclusion}

\begin{itemize}
  \item Limitations
  \item Perspectives
\end{itemize}

\chapterstar{Bibliographie}

\bibliographystyle{alpha}
\bibliography{typ}
\end{document}
